\section{Kết quả thực nghiệm}

% Slide 35
\begin{frame}[t,shrink=20]{Intent interpretation: kết quả tổng thể}
\scriptsize
\begin{center}
\scriptsize
\renewcommand{\arraystretch}{1.16}
\resizebox{\textwidth}{!}{%
\begin{tabular}{lrrrr}
\toprule
Method & Field F1 & Schema validity & Unsupported & Premature commit \\
\midrule
Raw template & 0.492 & 0.0\% & 0.0\% & 0.0\% \\
Schema only & 0.646 & 93.3\% & 93.6\% & 22.5\% \\
System + schema & 0.765 & 92.5\% & 37.6\% & 1.7\% \\
Few-shot + schema & 0.776 & 97.5\% & 38.9\% & 10.8\% \\
\textbf{Proposed} & \textbf{0.894} & \textbf{98.3\%} & \textbf{0.0\%} & \textbf{0.0\%} \\
\bottomrule
\end{tabular}%
}
\end{center}
\begin{alertblock}{Thông điệp}
deterministic validation tạo phần lớn safety gain.
\end{alertblock}
\end{frame}

% Slide 36
\begin{frame}[t,shrink=20]{Intent interpretation: các trường hợp còn khó}
\scriptsize
\vspace{0.15em}
\textbf{Kết quả theo nhóm}\par
\begin{itemize}
\item paraphrase: F1 0.907;
\item conflicting/ranked goals: 0.912;
\item indirect requests: 0.861;
\item compact syntax: 0.896.
\end{itemize}
\vspace{0.15em}
\textbf{Field sai nhiều}\par
\begin{itemize}
\item primary goal;
\item partition-change authorization;
\item GPU range;
\item resolution status;
\item allowed partitions.
\end{itemize}
\begin{alertblock}{Kết luận}
taxonomy và authorization semantics vẫn là bottleneck.
\end{alertblock}
\end{frame}

% Slide 37
\begin{frame}[t,shrink=20]{Profile feasibility và Pareto identification}
\scriptsize
\vspace{0.15em}
\textbf{Confirmation set}\par
\begin{itemize}
\item 6 unseen jobs;
\item 4 resource vectors;
\item 72 measurements;
\item median runtime CV 0.99\%.
\end{itemize}
\vspace{0.15em}
\textbf{Pareto identification}\par
\begin{itemize}
\item precision: 0.667;
\item recall: 0.847.
\end{itemize}
\vspace{0.15em}
\textbf{Raw recommendation}\par
\begin{itemize}
\item feasible rate: 83.3\%;
\item exact oracle match: 44.4\%.
\end{itemize}
\begin{alertblock}{Thông điệp}
exact match không phản ánh đầy đủ utility equivalence.
\end{alertblock}
\end{frame}

% Slide 38
\begin{frame}[t,shrink=20]{Safety–coverage trade-off}
\scriptsize
\begin{center}
\scriptsize
\renewcommand{\arraystretch}{1.16}
\resizebox{\textwidth}{!}{%
\begin{tabular}{lrrr}
\toprule
Policy & Commit rate & Feasible khi commit & Mean regret \\
\midrule
Original & 100.0\% & 66.7\% & 0.508 \\
Forced recommendation & 100.0\% & 83.3\% & 0.183 \\
\textbf{Probe-aware} & \textbf{38.9\%} & \textbf{100.0\%} & \textbf{0.020} \\
Oracle & 100.0\% & 100.0\% & 0.000 \\
\bottomrule
\end{tabular}%
}
\end{center}
\[
\text{Higher committed safety}
\Longleftrightarrow
\text{Lower automation coverage}
\]
\begin{exampleblock}{Hình minh họa}
commit rate và committed feasibility đặt cạnh nhau.
\end{exampleblock}
\end{frame}

% Slide 39
\begin{frame}[t,shrink=20]{Technical discontinuity tại biên 4 GPU}
\scriptsize
\begin{itemize}
\item 4 compatible jobs chạy thành công 12/12 repetitions;
\item 2 incompatible jobs fail 6/6;
\item failure xuất hiện ở preflight;
\item exit code 2;
\item không có compute result.
\end{itemize}
\begin{alertblock}{Kết luận}
feasibility phải được mô hình hóa bằng rule và preflight, không chỉ bằng predictor.
\end{alertblock}
\end{frame}

% Slide 40
\begin{frame}[t,shrink=20]{Prediction và uncertainty}
\scriptsize
\vspace{0.15em}
\textbf{Screening}\par
\begin{itemize}
\item factorized MAPE: 2.58\%;
\item kNN MAPE: 3.77\%;
\item monolithic MAPE: 19.54\%.
\end{itemize}
\vspace{0.15em}
\textbf{Confirmation}\par
\begin{itemize}
\item factorized MAPE: 21.04\%;
\item kNN MAPE: 19.13\%;
\item interval coverage: 94.12\%.
\end{itemize}
\vspace{0.15em}
\textbf{Interpretation}\par
Whole-job và unseen-resource holdout cho kết quả khó hơn nhưng đáng tin cậy hơn.\par
\end{frame}

% Slide 41
\begin{frame}[t,shrink=20]{Uncertainty thực sự thay đổi action}
\scriptsize
\begin{itemize}
\item interval coverage: 94.1\%;
\item schedule-probe rate: 61.1\%;
\item 18/18 decisions có đầy đủ trace;
\item 18/18 áp dụng đúng uncertainty threshold.
\end{itemize}
\begin{alertblock}{Thông điệp}
uncertainty không chỉ được báo cáo hậu nghiệm mà được chuyển thành action.
\end{alertblock}
\begin{exampleblock}{Hình minh họa}
\texttt{Interval narrow → commit}, \texttt{Interval wide → probe}.
\end{exampleblock}
\end{frame}

% Slide 42
\begin{frame}[t,shrink=20]{Queue-aware operational experiment}
\scriptsize
\vspace{0.15em}
\textbf{Cluster states}\par
\begin{itemize}
\item mostly free;
\item GPU fragmented;
\item controlled 2-GPU pressure.
\end{itemize}
\vspace{0.15em}
\textbf{Treatments}\par
\begin{itemize}
\item original request;
\item static Pareto;
\item queue-aware Pareto.
\end{itemize}
\vspace{0.15em}
\textbf{Correctness}\par
\begin{itemize}
\item 81/81 executions hoàn tất;
\item checksum đúng;
\item intent violation bằng 0.
\end{itemize}
\end{frame}

% Slide 43
\begin{frame}[t,shrink=20]{Operational utility và trade-off}
\scriptsize
\vspace{0.15em}
\textbf{Queue-aware so với original}\par
\begin{itemize}
\item fragmented: turnaround giảm trung bình 104.6 s;
\item controlled pressure: giảm 104.8 s;
\item mostly free: kém 11.1 s.
\end{itemize}
\vspace{0.15em}
\textbf{GPU-seconds}\par
Giảm trong mọi cluster state.\par
\vspace{0.15em}
\textbf{Makespan}\par
Không cải thiện phổ quát; dưới pressure có thể kém original khoảng 35–43 s.\par
\vspace{0.15em}
\textbf{Kết luận}\par
Queue-aware selection cải thiện user turnaround khi có pressure nhưng không đồng nghĩa với tối ưu batch makespan.\par
\end{frame}
